%!TEX root = Love602.tex
We may now regard \(n_1\) and \(\xi \) as known,
then equations (66) and (67) give \(t_1\) and
\(t_2\); and, \(\xi \) being \(dW/d\theta + U\),
the \nth{3} and \nth{4} of equations (65) give \(g_1\) and \(g_2\),
while \(n_2\) vanishes. Thus all the stress-resultants and
stress-couples are known.  Further, since
\[
\frac{dU}{d\theta} -U \cot(\theta)=\frac{a}{1-\sigma}(t_1-t_2)=
\frac{a}{1-\sigma}\left(\frac{d n_1}{d\theta} - n_1 \cot(\theta) \right),
\]
we have
\(U =\frac{a n_1}{1 - \sigma}+ b_1 \sin \theta \)
where \(b_1\) is a constant of integration, which may, however, be put equal to
zero, as the same term occurs in the extensional solution. Then the equation
\[
U \cot(\theta) -W= \frac{a (t_2 - \sigma t_1)}{1 - \sigma^2} b_1 \sin \theta
\]
gives \(W\).In order to satisfy given boundary-conditions of
force or displacement at the edges it is generally necessary to combine an
extensional solution with the edge-effect.

The equations determining the edge-effect in a spherical shell have not been
integrated, except in the case where n = 0. It may be observed here that, when
\(n=0\), the equations of equilibrium possess two integrals answering to
equations (63), and these can be shown to be
\[
n_1 \sin^2(\theta) -t_1 \sin (\theta) \cos (\theta) +
s_1 \sin (\theta) = \text{const},
\]

\[
\left(t_1 - \frac{g_1}{a} \right)\sin (\theta) -
\left(s_1 + \frac{h_1}{a} \right)\cos (\theta) = \text{const}.
\]

In order to determine the edge-effect in the case where \(n=1\), we should
begin by equating to zero the constants in the right-hand members of these equations.

